\documentclass[12pt]{article}
\usepackage[a4paper,vmargin=2cm,hmargin=3cm]{geometry}
\usepackage[utf8]{inputenc}
\usepackage{setspace}
\usepackage[english]{babel}
\setlength{\parindent}{0em}
\usepackage{lmodern} 
\usepackage{amsmath}
\usepackage{amssymb}
\usepackage{graphicx}
\usepackage[export]{adjustbox}
\onehalfspacing
\usepackage{scalerel,stackengine}
\usepackage{float}
\usepackage[dvipsnames]{xcolor}
\usepackage{amsmath}
\DeclareMathOperator*{\argmax}{arg\,max}
\DeclareMathOperator*{\argmin}{arg\,min}
\newcommand{\indep}{\perp \!\!\! \perp}

\title{\textbf{Syllabus}}
\author{\textbf{Econometrics (25117)}}
\date{Universitat Pompeu Fabra, 2024-2025}

\begin{document}
\maketitle

\section*{Course Details}
\begin{itemize}
    \item \textbf{Lecturer}\\
    Pierre Magontier
    \item \textbf{Sessions}\\
    18 Theory Sessions\\
    7 Seminar Sessions (starting on week 3)
    \item \textbf{Exams}\\
    Midterm Exam --- TBD\\
    Final Exam --- TBD
\end{itemize}
\section*{Course Description}
The course aims to introduce students to econometric methods and techniques. Econometrics lies at the intersection of economics and statistics, and econometric methods are widely used in applied work to estimate economic models, test hypotheses, and produce forecasts for future trends in relevant economic and financial variables.\\

The course is structured into two main blocks to provide a comprehensive understanding of econometric methods.

\begin{enumerate}
    \item Fundamental Inference Methods and OLS Regression --- In this initial block, students will learn essential statistical inference techniques that form the foundation of econometric analysis. They will be introduced to ordinary least squares (OLS) regression models, which are pivotal for estimating relationships between variables. This section will cover key concepts such as model specification, interpretation of coefficients, hypothesis testing, and the assumptions underlying OLS. Through practical examples and applications, students will gain hands-on experience in using OLS regression to analyze real-world data.
    \item Limitations of OLS and Alternative Models --- The second block addresses the limitations inherent in OLS regression models, particularly when dealing with more complex data structures or non-linear relationships. Students will explore various challenges such as multicollinearity, heteroscedasticity, and autocorrelation that can affect model accuracy. To enhance their analytical toolkit, the block will introduce alternative econometric models and data configurations.
\end{enumerate}

The teaching methodology is based on a combination of lectures and TA sessions. The lectures aim to deepen students' understanding of the topics presented. They introduce various econometric estimators, explore their theoretical properties, and provide illustrative examples and applications to highlight the empirical relevance of the theory.\\

The practical sessions are designed to equip students with problem-solving skills and data analysis capabilities. During these sessions, students will engage in exercises and are encouraged to participate actively. In these sessions, students will practice the techniques discussed in the lectures to estimate key parameters and conduct hypothesis tests from which they can draw meaningful conclusions.\\

Additionally, students are expected to dedicate time to self-study by reading the book chapter references as well as the supplementary materials suggested in class and working through practical questions to reinforce the skills learned during the lectures. This comprehensive approach ensures that students not only understand the theoretical aspects of basic econometrics but also become proficient in applying these methods in real-world scenarios.\\

By the end of the course, students should have an overall robust understanding of both foundational and basic econometric techniques, equipping them to tackle a wide range of economic and financial analysis challenges.


\section*{Grading}
To pass the course students have to obtain 50 out of 100 points according to the following distribution:
\begin{enumerate}
    \item Final exam: 60 points
    \item Mid-Term Exam: 30 points
    \item Seminars (submitting the solutions of exercise lists): 10 points
\end{enumerate}

Seminar exercises must be submitted individually to the teaching assistant at the beginning of each seminar. Printed copies or scanned versions of written solutions will not be accepted. Handwritten solutions are preferred to prevent any issues of academic integrity. Any copy-pasted responses will not be accepted, regardless of who contributed to the work.\\

There will be a second chance to take the final exam at the end of the quarter only  for those  who  failed  to  pass  the  course,  who  have  attended  seminars and  who  have submitted the solutions and summaries of the exercise lists. 


\section*{Requirements}
This course is intensive and requires a solid commitment to self-study. Students should be prepared to engage in readings, complete assignments, and participate in practical exercises. Active participation is encouraged, as it will enhance understanding and collaboration. A foundational background in probability, statistics, and economics is required for engaging with the course material. Students are expected to read the primary textbook references and any supplementary materials if need be.

\section*{Contents}

The core structure of the course follows \textit{Introduction to Econometrics, 4th Edition, Global Edition} by Stock and Watson. The selected chapters from this book are \underline{highly} recommended readings.

\begin{itemize}
    \item \textbf{Topic I: Introduction and Review of Probability}
    \begin{itemize}
        \item Introduction to Econometrics, 4th Edition, Global Edition, by Stock and Watson -- Chapter 1-2
        \item Introductory Econometrics, 5th Edition, A Modern Approach, by Jeff. Wooldridge -- Appendix B
        \item  Causal inference: The Mixtape, by Scott Cunningham (2021) -- Chapter 1-2

    \end{itemize}

    \item \textbf{Lecture II: Review of Statistics}
    \begin{itemize}
        \item Introduction to Econometrics, 4th Edition, Global Edition, by Stock and Watson -- Chapter 3
        \item Introductory Econometrics, 5th Edition, A Modern Approach, by Jeff. Wooldridge -- Appendix C\\[1em]
    \end{itemize}
    
    \item \textbf{Lecture III: Simple Linear Regression}
    \begin{itemize}
        \item Introduction to Econometrics, 4th Edition, Global Edition, by Stock and Watson -- Chapter 4
        \item Introductory Econometrics, 5th Edition, A Modern Approach, by Jeff. Wooldridge -- Chapter 2
    \end{itemize}
    
    \item \textbf{Lecture IV: OLS Implementation}
    \begin{itemize}
        \item Introduction to Econometrics, 4th Edition, Global Edition, by Stock and Watson -- Chapter 5
        \item Introductory Econometrics, 5th Edition, A Modern Approach, by Jeff. Wooldridge -- Chapter 8
    \end{itemize}
    
     \item \textbf{Lecture V: Multivariate Linear Regression}
     \begin{itemize}
        \item Introduction to Econometrics, 4th Edition, Global Edition, by Stock and Watson -- Chapter 6
        \item Introductory Econometrics, 5th Edition, A Modern Approach, by Jeff. Wooldridge -- Chapter 3
    \end{itemize}
    
     \item \textbf{Lecture VI: Inference and Interpretation}
     \begin{itemize}
        \item Introduction to Econometrics, 4th Edition, Global Edition, by Stock and Watson -- Chapter 7
        \item Introductory Econometrics, 5th Edition, A Modern Approach, by Jeff. Wooldridge -- Chapter 4-6
    \end{itemize}

    \item \textbf{Lecture VII: Functional Forms}
    \begin{itemize}
        \item Introduction to Econometrics, 4th Edition, Global Edition, by Stock and Watson -- Chapter 8
        \item Introductory Econometrics, 5th Edition, A Modern Approach, by Jeff. Wooldridge -- Chapter 6-9
    \end{itemize}

    \item \textbf{Lecture VIII: Threats to Identification}
    \begin{itemize}
        \item Introduction to Econometrics, 4th Edition, Global Edition, by Stock and Watson -- Chapter 1-2
        \item Introductory Econometrics, 5th Edition, A Modern Approach, by Jeff. Wooldridge -- Appendix B
        \item  Causal inference: The Mixtape, by Scott Cunningham (2021) -- Chapter 1-2

    \end{itemize}

    \item \textbf{Lecture IX: Functional Forms}
    \begin{itemize}
        \item Introduction to Econometrics, 4th Edition, Global Edition, by Stock and Watson -- Chapter 9
        \item  Causal inference: The Mixtape, by Scott Cunningham (2021) -- Chapter 3

    \end{itemize}

    \item \textbf{Lecture X: Instrumental Variables}
    \begin{itemize}
        \item Introduction to Econometrics, 4th Edition, Global Edition, by Stock and Watson -- Chapter 12
        \item Introductory Econometrics, 5th Edition, A Modern Approach, by Jeff. Wooldridge -- Chapter 15
        \item  Causal inference: The Mixtape, by Scott Cunningham (2021) -- Chapter 6
    \end{itemize}

    \item \textbf{Lecture XI: Panel Data}
    \begin{itemize}
        \item Introduction to Econometrics, 4th Edition, Global Edition, by Stock and Watson -- Chapter 10-13
        \item Introductory Econometrics, 5th Edition, A Modern Approach, by Jeff. Wooldridge -- Chapter 13-14
        \item  Causal inference: The Mixtape, by Scott Cunningham (2021) -- Chapter 7
    \end{itemize}

    \item \textbf{Lecture XII: Binary Outcomes}
    \begin{itemize}
        \item Introduction to Econometrics, 4th Edition, Global Edition, by Stock and Watson -- Chapter 11
        \item Introductory Econometrics, 5th Edition, A Modern Approach, by Jeff. Wooldridge -- Chapter 7-17
        \item  

    \end{itemize}

\end{itemize}




\end{document}